\section{Introduction}

    Understanding Biodiversity is one of the most important topic in biological studies. Why there are so many species, and how could they kept existing are the fundamental questions in biodiversity research. These questions were firstly considered as axiomatic and natural, since a stable ecosystem where there are many species is observed everyday. Biodiversity is viewed the intrinsic property of a ecological system, which could contribute to a stable function of the system. Furthermore, a huge mount of empirical studies indicate the point that high biodiversity represent a complex ecological system, and such a complicity makes the system stable, that is, could keep the state where many species coexist. However, due to the numerous and inconsistent definition of biodiversity as well as stability, there is still lack of concrete evidence and theory to prove the biodiversity-stability relation.
    
    In the contrast, ecosystem, especially community when considering the species existing and biodiversity, can be modeled as a dynamics system. With the increasing understanding of complex system as well as random matrix theory, May first argue that if the interaction strength between species is random and the system can be linearized around the stationary point, then large system is less probable to be stable. Such a statement based on the fact from random matrix theory that the distribution of eigenvalues of a large random matrix always converge to a round plane in complex panel, whose radius is related to the dimension of the system, that is, the number of species. Since in a continuous dynamics system, it is stable only if eigenvalues of the system is smaller than $1$, then the larger the system is, the less probable it is stable.

    As a result, stability is not a promising property of a large system, the species coexist should be considered from the specificity of ecosystem distinguishing from random system. One approach is to studying the structure of interactions, that is, the interaction  matrix of the system. Interaction types is important to the stability, which related to the signs of each interacting pairs. Studies has shown that predator-prey community is more probable to be stable compared to a mutualism or competing community with the same scale, because the eigenvalues of interaction matrix that symmetric elements has different sigh will distributed in a narrow ellipse. Besides, interaction network structure is important to the stability. The pattern of connecting, including nestedness, modularity, as well as network motifs and so on, indicate the complex connectivity and long-range relation in the community, thus being crucial to the stability. Empirical research has shown that in real community, nestedness normally negatively related to stability whereas modularity helps the maintain of coexistence. Another approach focus on the function of interaction as well as population growth. Species interaction as a complicated process in the nature, not always fit the simple GLV model well. Nonlinear reaction term, for example, Holling reaction functions, are introduced to the community dynamic models, and some of them has been proved to be much easier being stable with random matrices.

    However, besides of the dynamics of local community, spatial condition also contribute to the stable coexist of species. Island biogeography first predict that larger island, or patch, normally contain more species than the smaller ones. This hypothesis has been assured by multiple studies on different regions. Based on which, the neutral theory explained this result by assuming all species are neutral, that is, identical in dispersal and interaction, and gave a well prediction on the population size distribution. Furthermore, a recent study has shown that with a very large amount of patches, species could always coexist even only one can survive in each. On the contrast, large patch is more probable to contain more heterogeneous habitats, thus has more niche for different species. In the perspective of niche theory, high biodiversity in large space can be explained as the species turnover along the environment gradient. $\beta$-diversity is then used to measure the strength of species turnover. But in both sides, the spatial effect on biodiversity is related to the dispersal of individuals in the space, no matter randomly or purposely.

    To model the dispersal and interaction in space, the reaction-diffusion model is used. Firstly introduced to explaining the pattern formation during development, reaction-diffusion model is now well used in mathematical biology to simulate the result of spatial dispersal and interaction between particles, biological molecular, as well as individuals. Typically, in community ecology, reaction and diffusion in the model represent the interaction and dispersal of species. The traveling wave solution is mostly be considered, where there is a steep transition from one state, for example, some species, to another. A traveling wave solution with speed $0$ is also called a solitary, indicating the turnover of species without environment gradient, as well as species coexist in large scale, thus has important ecological meaning. One species case is equivalent to the condition where a species colonize a open area. Two species model also be well studied because of its complexity from interaction term. Though is the simplest case where there are interactions between species, the existence of traveling wave solution of the two species model is not fully solved. For system where there are multiple species, related problem getting even harder, considering the increase of dimensions when solving the whole system.

    Here, in this thesis project, we focus on the property of the potential traveling wave solution with speed $0$. With very simplified two species and multiple species model, we analyze and simulate the shape of the transition if it exists. Based on the population size distribution, we then proposed a $beta$-diversity measure, representing the species turnover in neutral environment.
